\documentclass[../../../include/open-logic-section]{subfiles}

\begin{document}
\olfileid{sth}{ord-arithmetic}{intro}

\olsection{سريزه}

په \olref[ordinals][]{chap} کښې مو د ترتيبي عددونو تيوري جوړه کړه۔
په \olref[spine][]{chap} کښې مو وليدل چې ترتيبي عددونه د هغې
ملاتير په توګه ګڼلے شو چې د سلسلې پاتې برخه پرې رغېږي۔ خو د
ترتيبي عددونو يوازينۍ ونډه دا نه ده۔ د ترتيبي عددونو حساب هم بايد وکړو۔

مخکښې په \olref[ordinals][idea]{sec} کښې مو د $\omega$،
$\omega+1$ او $\omega+\omega$ په يادولو دې خبرې ته اشاره کړې وه۔
هلته مو غير رسمي خبرې کړې وې؛ اوس ئې د کره بيانولو وخت دے۔ خو دا
هم ووايو چې په دې څپرکي کښې ډېر فلسفي بحث نشته؛ تر ډېره تخنيکي
پرمختګونه دي، او ورسره دا يو (لږ څه) په زړه پورې مشاهده چې يو کار
په دوو بېلو لارو کولے شو۔

\end{document}
